\clearpage
\phantomsection\label{references}
\addcontentsline{toc}{section}{References and Source Roles}
\renewcommand{\refname}{References and Source Roles}
\begin{thebibliography}{99}
\bibitem{xiang2024}
\emph{Giant magnetocaloric effect in spin supersolid candidate Na$_2$BaCo(PO$_4$)$_2$}, Nature \textbf{625}, 270--275 (2024).
\url{https://doi.org/10.1038/s41586-023-06885-w}

\bibitem{nmr2025}
\emph{NMR study of spin supersolid phases in the triangular-lattice antiferromagnet Na$_2$BaCo(PO$_4$)$_2$}, Physical Review B \textbf{112}, 125163 (2025).
\url{https://doi.org/10.1103/gsk8-1k9q}

\bibitem{kcso2026}
\emph{Phase diagram and spectroscopic signatures of a supersolid in the quantum Ising magnet K$_2$Co(SeO$_3$)$_2$}, Nature Communications (2026).
\url{https://doi.org/10.1038/s41467-026-69661-0}

\bibitem{zhu2026}
M. Zhu et al., \emph{Transverse order and longitudinal fluctuations in a near-Ising spin supersolid}, arXiv:2607.13635v2 (2026), preprint.
\url{https://arxiv.org/abs/2607.13635v2}

\bibitem{rb2026}
\emph{Spin-supersolidity induced quantum criticality and magnetocaloric effect in the triangular-lattice antiferromagnet Rb$_2$Co(SeO$_3$)$_2$}, npj Quantum Materials (2026).
\url{https://doi.org/10.1038/s41535-026-00881-9}

\bibitem{gao2022}
Y. Gao et al., \emph{Spin supersolidity in nearly ideal easy-axis triangular quantum antiferromagnet Na$_2$BaCo(PO$_4$)$_2$}, npj Quantum Materials \textbf{7}, 89 (2022).
\url{https://doi.org/10.1038/s41535-022-00500-3}

\bibitem{park2026}
S. Park, S.-M. Park, Y.-T. Oh, H.-Y. Lee and E.-G. Moon, \emph{Spin-orbit-induced Instability and Finite-Temperature Stabilization of a Triangular-lattice Supersolid}, arXiv:2601.20963v1 (2026). Source for the model, common-$z$ orbit, Figure 4 parameters, Section II and Eqs. (3)--(4).
\url{https://arxiv.org/abs/2601.20963v1}

\bibitem{chen2017}
J. Chen et al., \emph{Phase transition of the q-state clock model: duality and tensor renormalization}, Chinese Physics Letters \textbf{34}, 050503 (2017). The abstract reports tensor calculations of both transition temperatures for the standard six-state clock model; this supports the effective-model mechanism, not microscopic NBCP matching.
\url{https://arxiv.org/abs/1706.03455}

\bibitem{lin2025}
X. Lin and T. Shi, \emph{Pseudo-Goldstone Modes at Finite Temperature}, arXiv:2505.07229v1 (2025). Reference for nonlinear soft paths; the internal zero-field XXZ mode is a separate target.
\url{https://arxiv.org/html/2505.07229v1}

\bibitem{rau2018}
J. G. Rau, P. A. McClarty and R. Moessner, \emph{Pseudo-Goldstone gaps and order-by-quantum-disorder in frustrated magnets}, Physical Review Letters \textbf{121}, 237201 (2018). Source for the leading-order curvature formula and conjugate zero-mode coordinates; the susceptibility reduction applies it to the common-$z$ mode.
\href{https://doi.org/10.1103/PhysRevLett.121.237201}{doi:10.1103/PhysRevLett.121.237201};
\url{https://arxiv.org/html/1805.00947v2}

\bibitem{jose1977}
J. V. Jos\'{e}, L. P. Kadanoff, S. Kirkpatrick and D. R. Nelson, \emph{Renormalization, vortices, and symmetry-breaking perturbations in the two-dimensional planar model}, Physical Review B \textbf{16}, 1217 (1977); erratum, \textbf{17}, 1477 (1978), \href{https://doi.org/10.1103/PhysRevB.17.1477}{doi:10.1103/PhysRevB.17.1477}. The publisher abstract supports the sixfold-anisotropy mechanism; the full article and erratum were not inspected. The linearized exponents are derived explicitly in this note.
\url{https://doi.org/10.1103/PhysRevB.16.1217}

\end{thebibliography}

\phantomsection\label{sources-for-the-gap-analysis}
Gap sources: \cite{park2026,rau2018,lin2025}.
\phantomsection\label{sources-for-the-clock-analysis}
Clock sources: \cite{chen2017,park2026,jose1977}.
